\subsection{\texorpdfstring{ISOMORPHISMS BETWEEN \(\wedge^{p}(\mathbf{M})\) AND \(\wedge^{n - p}(\mathbf{M}^{*})\) FOR AN \(n\) -DIMENSIONAL FREE MODULE M}{ISOMORPHISMS BETWEEN \textbackslash wedge\^{}\{p\}(\textbackslash mathbf\{M\}) AND \textbackslash wedge\^{}\{n - p\}(\textbackslash mathbf\{M\}\^{}\{*\}) FOR AN n -DIMENSIONAL FREE MODULE M}}

PROPOSITION 12. Let \(\mathbf{M}\) be a free A-module of dimension \(n\) ; let \(e \in \bigwedge^{n}(\mathbf{M})\) be an element forming a basis of \(\bigwedge^{n}(\mathbf{M})\) and let \(e^{*}\) be the element of \(\bigwedge^{n}(\mathbf{M}^{*})\) such that \(\{(-1)^{n(n-1)/2} \theta_{\bigwedge}(e^{*})\}\) is the dual basis of \(\{e\}\) in \((\bigwedge^{n}(\mathbf{M}))^{*}\) . Let \(\phi: \bigwedge(\mathbf{M}^{*}) \to \bigwedge(\mathbf{M})\) be the mapping \(z \mapsto z \lrcorner e^{*}\) and \(\phi': \bigwedge(\mathbf{M}^{*}) \to \bigwedge(\mathbf{M})\) the mapping \(z^{*} \mapsto z^{*} \lrcorner e\) . Let \(\phi_{p}\) (resp. \(\phi_{p}'\) ) be the restriction of \(\phi\) (resp. \(\phi'\) ) to \(\bigwedge^{p}(\mathbf{M})\) (resp. \(\bigwedge^{p}(\mathbf{M}^{*})\) ). Then:

\begin{enumerate}
\def\labelenumi{(\roman{enumi})}
\item
  The mapping \(\phi\) is a left \(\bigwedge (\mathbf{M})\) -module isomorphism and the mapping \(\phi'\) is a left \(\bigwedge (\mathbf{M}^*)\) -module isomorphism; moreover the mappings \(\phi\) and \(\phi'\) are inverses of one another.
\item
  The mapping \(\phi_p\) is an isomorphism of the A-module \(\bigwedge^p (\mathbf{M})\) onto the A-module \(\bigwedge^{n - p}(\mathbf{M}^*)\) and the mapping \(\phi_p'\) is an isomorphism of the A-module \(\bigwedge^p (\mathbf{M}^*)\) onto the A-module \(\bigwedge^{n - p}(\mathbf{M})\) .
\item
  If we write \(\mathbf{B}(u,v^{*}) = \langle u,\theta_{\wedge}(v^{*})\rangle\) for \(u\in \bigwedge (\mathbf{M})\) and \(v^{*}\in \bigwedge (\mathbf{M}^{*})\) then, for \(u^{*}\in \bigwedge^{p}(\mathbf{M}^{*})\) and \(v^{*}\in \bigwedge^{n - p}(\mathbf{M}^{*})\)
\end{enumerate}

\[
\mathbf {B} \left(\phi_ {p} ^ {\prime} (u ^ {*}), v ^ {*}\right) = (- 1) ^ {p (n - p)} \mathbf {B} \left(u ^ {*}, \phi_ {n - p} ^ {\prime} (v ^ {*})\right).\tag{77}
\]

The fact that \(\phi\) is \(\bigwedge (\mathbf{M})\) -linear and \(\phi'\) is \(\bigwedge (\mathbf{M}^*)\) -linear follows from the formulae \((u \wedge v) \lrcorner e^* = u \lrcorner (v \lrcorner e^*)\) and \((u^* \wedge v^*) \lrcorner e = u^* \lrcorner (v^* \lrcorner e)\) (no. 6, formula (37), using the fact that \(\theta_{\wedge}\) is an isomorphism of \(\bigwedge (\mathbf{M}^*)\) onto the opposite algebra to \(\bigwedge (\mathbf{M}^*)\) . On the other hand there exists a basis \((e_i)_{1 \leq i \leq n}\) of \(\mathbf{M}\) such that

\[
e = e _ {1} \wedge e _ {2} \wedge \dots \wedge e _ {n} \quad \text { and } \quad e ^ {*} = (- 1) ^ {n (n - 1) / 2} e _ {1} ^ {*} \wedge e _ {2} ^ {*} \wedge \dots \wedge e _ {n} ^ {*},
\]

where \((e_i^*)\) is the basis dual to \((e_i)\) . We write \(\mathbf{I} = [1, n]\) ; it follows from (76) that, for every subset \(\mathbf{J}\) of \(\mathbf{I}\) with \(p\) elements,

\[
\left\{ \begin{array}{l} \phi (e _ {J}) = (- 1) ^ {n (n - 1) / 2 + p (p - 1) / 2} \rho_ {J, I - J} e _ {I - J} ^ {*} \\ \phi^ {\prime} (e _ {J} ^ {*}) = (- 1) ^ {p (p - 1) / 2} \rho_ {J, I - J} e _ {I - J}. \end{array} \right.\tag{78}
\]

This proves that \(\phi\) and \(\phi'\) are bijective; moreover \(\rho_{J,I-J}\rho_{I-J,J}=(-1)^{p(n-p)}\) ( \(\S 7\) , no. 8, formula (21)); as the number

\[
\frac {n (n - 1)}{2} + \frac {p (p - 1)}{2} + \frac {(n - p) (n - p - 1)}{2} + p (n - p) = n (n - 1)
\]

is even, it follows that \(\phi\) and \(\phi'\) are inverses of one another. Finally, to prove (77), it suffices to take \(u^{*} = e_{J}^{*}\) and \(v^{*} = e_{I - J}^{*}\) ; the verification also follows from the definition of \(\theta_{\wedge}\) , formulae (78) and the relation \(\rho_{J,I - J}\rho_{I - J,J} = (-1)^{p(n - p)}\) (§ 7, no. 8, formula (21)). Note that, for \(u^{*}\in \bigwedge^{p}(\mathbf{M}^{*})\) and \(v^{*}\in \bigwedge^{n - p}(\mathbf{M}^{*})\) , \(\mathbf{B}(\phi_{p}^{\prime}(u^{*}), v^{*})\) is, to within a sign, the coefficient of \(u^{*} \wedge v^{*}\) with respect to the basis \(\{e^{*}\}\) of \(\bigwedge^{n}(\mathbf{M}^{*})\) .

PROPOSITION 13. With the hypotheses and notation of Proposition 11, for every endomorphism g of the A-module M.

\[
(\det g) \phi = \bigwedge (^ {t} g) \circ \phi \circ \bigwedge (g).\tag{79}
\]

Clearly \(\bigwedge(^{t}g) = \theta_{\wedge}^{-1}\circ (^{t}\bigwedge (g))\circ \theta_{\wedge}\) ; since \(\bigwedge (g)\) is an endomorphism of the algebra \(\bigwedge (\mathbf{M})\) and by definition, for all

\[
z \in \bigwedge (\mathbf {M}), \quad \theta_ {\wedge} (\bigwedge (g) (z) \lrcorner e ^ {*}) = \theta_ {\wedge} (e ^ {*}) \llcorner \theta_ {\wedge} (\bigwedge (g) (z)),
\]

we deduce from formula (42) of no. 6 that

\[
\begin{array}{r l} ((\theta_ {\wedge} ^ {- 1} \circ (^ {t} \bigwedge (g)) \circ \theta_ {\wedge}) \circ \phi \circ \bigwedge (g)) (z) & = \theta_ {\wedge} ^ {- 1} (^ {t} \bigwedge (g) (\theta_ {\wedge} (e ^ {*})) \llcorner z) \\ & = z \lrcorner (\bigwedge (t g) (e ^ {*})) = (\det g) (z \lrcorner e ^ {*}) = (\det g) \phi (z) \end{array}
\]

taking account of § 8, no. 4, Proposition 8.

COROLLARY. For every automorphism g of E,

\[
\bigwedge(^ {t} g ^ {- 1}) = (\det g) ^ {- 1} \phi \circ (\bigwedge (g)) \circ \phi^ {- 1}.\tag{80}
\]

